\chapter*{Abstract}
\addcontentsline{toc}{chapter}{Abstract}

Academic publication lists make research visible but do not, by themselves, support full-text discovery, evidence-linked question answering, or the translation of prior work into feasible student projects. This thesis presents an evidence-based case study of the TTLAB Research Intelligence Platform, a local research-discovery prototype built with FastAPI, SQLite, React, and deterministic or locally hosted language-model components. The audited snapshot contains 134 publication records, 98 extracted PDFs, 698 pages, and 756 page-aware text chunks. It implements keyword, hashing-vector, and hybrid retrieval; citation-returning question answering; thesis-extension recommendations; topic and author exploration; generated paper artefacts; and an administrative review mechanism.

The contribution is a reproducible, source-traceable engineering artefact and a rigorous characterization of its current limits, rather than a claim of a novel retrieval algorithm or demonstrated user benefit. Repository inspection, database integrity checks, isolated automated tests, API probes, frontend compilation, and interface inspection establish implementation evidence. The active hashing-vector index covers only 25 of 756 chunks, 32.94\% of chunks have an unknown inferred section, author identity defects remain, and no completed gold retrieval, question-answering, or human-review study exists. Consequently, persisted ``grounded'' statuses are interpreted only as structural and lexical citation checks, not factual accuracy. The thesis connects each research question to a method, evidence source, result, and bounded conclusion, and specifies the evaluation work required before effectiveness claims can be made.

\textbf{Keywords:} research intelligence, information retrieval, retrieval-augmented generation, scholarly recommendation, citation grounding, research software, reproducibility.
